% ============================================================================
%  Computer Science for Physics and Chemistry -- GROUP RESEARCH PAPER TEMPLATE
%  ONE single file. Write directly into it. Replace the guidance text in each
%  section with your own writing. Compiles on Overleaf as-is.
%  >>> DELETE the grey "Guidance:" lines and the Instructions box before you
%      submit the final version. <<<
% ============================================================================

\documentclass[sigconf,nonacm]{acmart}
% 'nonacm' removes ACM copyright/CCS boilerplate. If acmart is unavailable, see
% the fallback note at the end of this file.

\settopmatter{printacmref=false}
\renewcommand\footnotetextcopyrightpermission[1]{}
\pagestyle{plain}

\usepackage{graphicx}
\usepackage{booktabs}
\usepackage{xcolor}
\usepackage{tcolorbox}

% a grey "Guidance" style for the inline hints (delete these before submitting)
\definecolor{guide}{RGB}{110,120,130}
\newcommand{\guidance}[1]{\par{\small\color{guide}\textit{Guidance: #1}}\par\medskip}

\begin{document}

% ===========================================================================
%  INSTRUCTIONS BOX  (visible in the PDF; DELETE before final submission)
% ===========================================================================
\begin{tcolorbox}[colback=yellow!8,colframe=orange!60!black,title=\textbf{READ FIRST --- then delete this box before submitting}]
\small
\textbf{What this is.} A short scientific paper (4--6 pages, two columns)
reporting a small study \emph{you} carry out: pick a problem, find a public
dataset (\textbf{UCI} or \textbf{Kaggle}), analyse it in Python with the course
methods, get real numbers and plots, and discuss them.

\smallskip
\textbf{Three hard requirements --- miss any one and the paper is NOT accepted:}
(1)~a \textbf{GitHub repository} link; (2)~the \textbf{dataset} inside the repo
\emph{or} a link in \textbf{both} the repo README \textbf{and} this paper;
(3)~a \textbf{detailed AI-usage disclosure}.

\smallskip
\textbf{Grading (out of 20).}
\textit{Structure \& presentation} --- 4
(sections present + correct two-column formatting + clear writing).
\textit{Paper content} --- 9
(problem \& cited data 1.5; methods 2.5; results 2.5; discussion 2.0;
abstract \& conclusion 0.5).
\textit{Source code \& reproducibility} --- 7
(repo builds \& runs 2.5; reproduces the paper 2.0; data handling 1.0;
code hygiene 1.0; AI disclosure 0.5).

\smallskip
The best papers are invited to the \textbf{poster session}, competing for a prize.
\end{tcolorbox}

% ---------------------------------------------------------------------------
%  TITLE AND AUTHORS
% ---------------------------------------------------------------------------
\title{Your Paper Title: Short, Specific, Descriptive}

\author{First Author}
\affiliation{\institution{UFAZ}\country{Azerbaijan}}
\email{name@example.com}

\author{Second Author}
\affiliation{\institution{UFAZ}\country{Azerbaijan}}
\email{name@example.com}
% Add or remove \author blocks to match your group size.

\begin{abstract}
\guidance{The abstract is one self-contained paragraph of about 150 words that
a reader can understand without seeing the rest of the paper. In one or two
sentences each, state: the problem or question you studied; the dataset you
used; what you did to it (your method); your single most important result, as a
concrete finding or number; and one sentence on what it means. Do not put
citations, figures, or equations in the abstract. Write it last, once the paper
is finished.}
Write your abstract here.
\end{abstract}

\maketitle

% ---------------------------------------------------------------------------
\section{Introduction}
\guidance{Open by saying what problem or question you are tackling and why it
matters to someone. Give the reader just enough background to care, but do not
write a long literature review. Briefly say what data you will use (one line;
the full description goes in the Data section). State clearly, in one or two
sentences, what you set out to do, beginning with ``In this work, we\ldots''.
If you like, end with a one-sentence road map of the paper.}
Write the introduction here.

% ---------------------------------------------------------------------------
\section{Data}
\guidance{Name your dataset and the public source it comes from (UCI or Kaggle)
and cite it with \texttt{\textbackslash cite\{dataset\}}~\cite{dataset}. Give
the link, and remember the \emph{same} link must appear in your repository
README. Describe the dataset concretely: how many samples (rows), what the
variables/columns are, and their units. Explain any cleaning or preprocessing
you did and why (removing missing values, filtering, scaling). Mention anything
unusual such as outliers, noise, or imbalance.}
Describe your dataset, its source, size, and any preprocessing here.

% ---------------------------------------------------------------------------
\section{Methods}
\guidance{Explain your computational approach clearly enough that another
student could reproduce it. State which course methods you used (for example
distribution fitting, Q--Q plots and the Shapiro--Wilk test, a p-value test,
correlation, curve fitting, or Monte~Carlo) and why each was appropriate for
your question. Name your tools: Python with NumPy, SciPy, and Matplotlib, run
inside your conda environment, and note that the exact versions are pinned in
\texttt{environment.yml}. If you fit or test something, say exactly what you
compared and which function you called. Do not paste large blocks of code here;
the code lives in your repository.}
Describe your methods and tools here.

% ---------------------------------------------------------------------------
\section{Results}
\guidance{Report what you actually found, using real numbers and plots produced
by your own code. Introduce each figure and table in the text and refer to it by
number, for example ``Figure~\ref{fig:example} shows\ldots''. Every figure and
table must have a caption that explains what it shows. Where it makes sense,
report uncertainty: a value together with its error, or a p-value for a
comparison. State the results plainly and honestly, including any that did not
turn out as expected, and save the interpretation for the Discussion.}

\begin{figure}[t]
  \centering
  \includegraphics[width=\linewidth]{example-image}
  % replace example-image with your own plot file
  \caption{What the figure shows, in one sentence.}
  \label{fig:example}
\end{figure}

\begin{table}[t]
  \centering
  \caption{What the table reports.}
  \label{tab:example}
  \begin{tabular}{@{}lr@{}}
    \toprule
    Quantity & Value \\
    \midrule
    mean     & 0.00 \\
    p-value  & 0.00 \\
    \bottomrule
  \end{tabular}
\end{table}

Write your results here, referring to Figure~\ref{fig:example} and
Table~\ref{tab:example}.

% ---------------------------------------------------------------------------
\section{Discussion}
\guidance{This is where you interpret your results, and it carries real marks,
so do not skip it. Answer the question you raised in the Introduction: what do
the numbers actually mean? Be honest about the limitations of your study, such
as a small or noisy sample, or assumptions that may not hold (for instance, did
you assume the data was normal, and was it really?). Say what you would do next
with more time or more data. Remember that ``it worked'' is not a discussion:
explain \emph{why} your results are what they are.}
Interpret your results, state limitations, and suggest next steps here.

% ---------------------------------------------------------------------------
\section{Conclusion}
\guidance{Write two or three sentences only. Remind the reader what you set out
to do, and state the single most important thing you found. Do not introduce any
new results or ideas here.}
Write a short conclusion here.

% ---------------------------------------------------------------------------
\section{Reproducibility}
\guidance{This section is required: with no working repository the paper is not
accepted. Give the full link to your GitHub repository. State that the
repository contains your code, the \texttt{environment.yml}, a README with run
instructions, and the data (or the data link). Give the one or two commands that
reproduce your results, for example: \texttt{conda env create -f
environment.yml}, then \texttt{conda activate <env>}, then \texttt{python
main.py}.}
\textbf{Repository:} \url{https://github.com/USER/REPO}

State what the repository contains and how to run it here.

% ---------------------------------------------------------------------------
\section{AI Usage Disclosure}
\guidance{This section is required and must be detailed. List every AI tool you
used (for example ChatGPT, Claude, GitHub Copilot) and state exactly what you
used each one for: writing or editing text, generating or debugging code,
analysing data, or making figures. Be specific about which parts of the work
each tool touched. If you used no AI tools at all, say so explicitly. Vague or
missing disclosure is treated as an academic-integrity issue.}
Disclose, in detail, any AI tools you used and what for, or state that none were
used.

% ---------------------------------------------------------------------------
%  REFERENCES
% ---------------------------------------------------------------------------
\begin{thebibliography}{9}

\bibitem{dataset}
Dataset Provider. Year.
\textit{Dataset Name}.
UCI Machine Learning Repository / Kaggle.
\url{https://dataset-link}

\bibitem{scipy}
Pauli Virtanen et al. 2020.
SciPy 1.0: Fundamental Algorithms for Scientific Computing in Python.
\textit{Nature Methods} 17, 261--272.

\end{thebibliography}

\end{document}

% ============================================================================
%  FALLBACK IF 'acmart' IS NOT AVAILABLE
%  Overleaf has acmart built in (New Project -> Templates -> "ACM Conference").
%  If you ever need a fallback, replace the \documentclass line with
%     \documentclass[twocolumn,10pt]{article}
%  and delete the \affiliation, \email, and \settopmatter lines. All section
%  content and structure stay the same.
% ============================================================================